\documentclass{xepersian-thesis-iasbs-ftex}
\usepackage[utf8]{inputenc}
\usepackage[T1]{fontenc}
\usepackage{geometry}
\usepackage{graphicx}
\usepackage{float}
\usepackage{hyperref}
\usepackage[english,indonesian]{babel}

\title{Optimasi Pemrosesan Template LaTeX Menggunakan Pendekatan Berbasis AI}
\author{Academic Support Agent}
\date{\today}

\begin{document}
\maketitle

\begin{abstract}

\end{abstract}

\chapter{PENDAHULUAN}


\chapter{TINJAUAN PUSTAKA}


% Original content (commented out):
% % -------------------- DON'T EDIT ----------------------------------------
\documentclass[a4paper,12pt,oneside,openany]{xepersian-thesis-iasbs-ftex}
\input{definitions}

% ------------------- You Can EDIT ----------------------------------------
\graphicspath{{chap1_images/}{chap2_images/}}                       % The paths to the images.

% -------------------- DON'T EDIT ----------------------------------------

\begin{document}

% -------------------- PLEASE EDIT ---------------------------------------

% --------------------------------------   INFORMATION IN PERSIAN  ---------------------------------------

\title{پایان‌نامه من با زی‌پرشین}
\author{نام و نام خانوادگی}
%--------------------------------------------
% از بین دو حالت زیر یکی را فعال کنید:
\degree{کارشناسی ارشد}
%\degree{دکتری}
%---------------------------------------------

% از بین دو حالت زیر یکی را فعال کنید:
\supervisor{نام استاد راهنما} \singlesupervisortrue
% \supervisors{نام استاد اول}{نام استاد دوم}\multisupervisorstrue
% در صورتی که استاد مشاور داشته‌اید نام وی را در خط زیر بنویسید در غیر این صورت خط زیر باید غیرفعال باشد.
% \advisor{نام استاد مشاور} \advisorexisttrue
\department{علوم رایانه و فناوری اطلاعات}
\university{دانشگاه تحصیلات تکمیلی در علوم پایه زنجان}
\city{زنجان}
\thesisdate{\today}
\makepersiantitle
% -------------------------------------------------------------------------------------------------------------------
\startpage
% -------------------------------------------------------------------------------------------------------------------

\beginto
%%در اینجا پایان‌نامه خود را به هر کس که دوست دارید تقدیم کنید. در غیر این صورت این قسمت را غیر فعال کنید.
آنهایی که می‌خواهند بیشتر بدانند.
\endto
% -------------------------------------------------------------------------------------------------------------------
%اگر در فایل thanks از دوستان و همکاران خود تشکر کرده‌اید خط زیر را فعال کنید، در غیر این صورت این خط باید غیرفعال باشد.
\thanks
% -------------------------------------------------------------------------------------------------------------------
\newpage
\begin{abstract}
\addcontentsline{toc}{section}{چکیده} % To add the abstract in the index.
\input{abstract_fa}
\end{abstract}
% -------------------------------------------------------------------------------------------------------------------

\tableofcontents
\listoffigures
%\listoftables
% -------------------------------------------------------------------------------------------------------------------

% اگر قبل از شروع فصل‌های پایان‌نامه مقدمه‌ای دارید، خط زیر را فعال کرده و مقدمه خود را در فایل introduction بنویسید. در غیر این صورت این خط باید غیر فعال باشد.
\newpage
\pagenumbering{arabic}\setcounter{page}{1}
\include{introduction}

% ----------------------------------------   Chapters of the Thesis  ---------------------------------------------
\include{chapter_1} % if any un-wanted empty page is produced between chapters, you can use \input instead of \include : \input{chapter_1}
\include{chapter_2}
\include{chapter_3}
\include{chapter_4}

%--------------------------------------------------------------------------------------------------------
% -------------------- DON'T EDIT ----------------------------------------
% the following lines are needed for making the appendixes name correct in the index.
\makeatletter
\def\@makechapterhead#1{%
  \vspace*{50\p@}%
  {\parindent \z@ \centering\normalfont
    \ifnum \c@secnumdepth >\m@ne
      \if@mainmatter
        \huge\bfseries \@chapapp\space \thechapter
        \par\nobreak
        \vskip 20\p@
      \fi
    \fi
    \interlinepenalty\@M
    \Huge \bfseries #1\par\nobreak
    \vskip 40\p@  }}
%--------------------------------------------------------------------------------------------------------
% -------------------- Please EDIT ----------------------------------------

\appendix
\include{appendix_1} 
\include{appendix_2}

\bibliographystyle{unsrt-fa}
% اگر فایل bibtex با پسوند bib حاوی اطلاعات مربوط به مراجع خود با فرمت صحیح bibtex را دارید از خط زیر استفاده کنید و به جای MyReferences نامه فایل خود را بنویسید.

% در این نمونه پایان‌نامه فرض شده است که شما فایل bib حاوی اطلاعات مربوط به مراجع خود با نام MyReferences.bib را دارید.
\bibliography{MyReferences}

% اگر به صورت عادی می‌خواهید ارجاع دهید خط بالا را غیر فعال کرده و قسمت زیر را فعال کنید و طبق مثال عمل کنید (البته این روش حرفه‌ای نیست و توصیه نمی‌شود).
%\begin{thebibliography}{99}
%\end{thebibliography}

\include{FE} % include persian to english dictionary

% --------------------------------------   INFORMATION IN LATIN  ----------------------------------------------
\begin{latin}
\latintitle{My Thesis in XePersian}
\latinauthor{First Name Last Name}
% --------------------------------------
% choose and activate one of the following lines
\latindegree{Master's Thesis}
%\latindegree{Ph.D. Thesis}
% --------------------------------------
\latinthesisdate{\latintoday}
% If you have an advisor, write its name in the following line, otherwise inactive (comment) the line.
\latinsupervisor{Name of Supervisor} \singlesupervisortrue
% \latinsupervisors{Name of the first supervisor}{Name of the second supervisor} \multisupervisorstrue
% If you have an advisor, write its name in the following line, otherwise inactive (comment) the line.
\latinadvisor{Name of Advisor} \advisorexisttrue
\latindepartment{Physics}
\latinuniversity{Institute for Advanced Studies in Basic Sciences}
\latincity{Zanjan}
\begin{latinabstract}
\input{abstract_en}
\end{latinabstract}
\makelatintitle
\end{latin}

\end{document}


\end{document}